\scanpage{229}

over such an local ring:

\uline{Theorem} Let $A$ be an analytic algebra over $\underline{C}$ (we can suppose $A$ to be the ring of convergent power series\deleted{\unsure{)}} in $n$ variables), $\widehat{A}$ its completion, $Y$ and $\widehat{Y}$ the spectra, $X$ a proper scheme over $Y$, $\widehat{X}=X\times_Y\widehat{Y}$. Then:

(i) The inverse system $(\operatorname{Pic}(X_n))$ satisfies MLAR (as stated above, this depends only on the car. 0 assumption for the residue field).

(ii) $\operatorname{Pic}(X)$ and $\operatorname{Pic}(\widehat{X})$ have same image in $\unsure{\operatorname{Pic}}(X_n)$, - namely the group of ``universal images''. (NB recall $\operatorname{Pic}(\widehat{X})\simeq\varprojlim\operatorname{Pic}(X_n)$).

(iii) In order for $\operatorname{Pic}(X)\longrightarrow\operatorname{Pic}(\widehat{X})$ to be an isomorphism, it is necessary and sufficient that $\operatorname{supp}R^1f_*(\underline{O}_X^{\deleted{\unsure{\bullet}}})\subset(y)$, i.e. $H^1(X,\underline{O}_X)$ of finite length over $A$. (NB It amounts also to the same to \deleted{\unsure{impo}} ask that the inverse system of the subgroups $\operatorname{Pic}'(X_n)$ of universal images is constant for large $n$).

We get for example:

\uline{Corollary} 1 The following conditions are equivalent: (i) $X/Y$ projective (ii) $\widehat{X}/\widehat{Y}$ projective (iii) For every $n$, $X_n/Y_n$ projective.

This applies for instance if $\dim X_0\leq1$.

Now, applying your theorem of \uline{resolution of singularities}, and Mumford's method of relating the local Picard group of $A$ to the global Picard group of a regular scheme dominating $A$ birationally, one gets from the last statement \deleted{\unsure{in}} (iii) of last theorem:

\uline{Corollary} 2. Let $A$ be as above, assume \deleted{\unsure{the closed point $y$ of $Y$ to}} $Y'=Y-(y)$ regular, and let $\widehat{Y}'=\widehat{Y}-(\widehat{y})$. Then $\operatorname{Pic}(Y')\longrightarrow\operatorname{Pic}(\widehat{Y}')$ is an isomorphism\deleted{\unsure{?}}.

This explains ``à priori'' (when $A$ is normal) why Mumford was

\margin{228}
